% Listado de requisitos funcionales, requisitos de datos y restricciones semánticas.

% Funcionales: 
% \begin{itemize}
%     \item Crear o mencionar hashtag
%     \item Listar tendencias
%     \item Asignar una categoría a un hashtag
%     \item Mostrar categoría ordenada
%     \item Eliminar tendencia
% \end{itemize}

%======================= %
% REQUISITOS FUNCIONALES %
% ======================= %


\section{Crear o mencionar hashtag}
\textbf{\textcolor{rojo}{\underline{RF2.1:}}} Crear o mencionar hashtag. \\[6pt]

\textbf{\textcolor{rojo}{\underline{Descripción:}}} Cada vez que haya una publicación con una cadena que empiece por "\#", almacenarlo de cara a poder ver los temas que son tendencia. En caso de ya existir, se actualiza el contador de veces que ha aparecido ese hashtag; si no, se crea y se pone el contador a uno.
 \\[6pt]

\textbf{\textcolor{rojo}{Entrada:}} Agente externo: usuario. Acción: realizar una publicación que contenga una cadena que empiece por "\#". 
Requisito de datos de entrada \textcolor{azul}{RDE2.1}. \\[6pt]

\textbf{\textcolor{rojo}{BD:}} Requisito de datos de escritura \textcolor{azul}{RDW2.1} y lectura \textcolor{azul}{RDR2.1.} \\[6pt]

\textbf{\textcolor{rojo}{Salida:}} Ninguna \\[6pt]

\textcolor{azul}{\textbf{RDE2.1:}} Datos de entrada de crear o mencionar hashtag \\
\hspace*{1cm} Hashtag: Cadena de caracteres (32) \\[12pt]

\textcolor{azul}{\textbf{RDR2.1:}} Datos para comprobar si ya existía un hashtag. \\
\hspace*{1cm} Los mismos que RDE2.1\\

\textcolor{azul}{\textbf{RDW2.1:}} Datos almacenados/actualizados de un hashtag. \\
\hspace*{1cm} Los mismos que RDE2.1 además de el contador de veces que ha aparecido ese hashtag.\\
\hspace*{1cm} Contador: Entero \\[12pt]

\textbf{\textcolor{verde}{Restricciones Semánticas:}} \\[6pt]
\textbf{\textcolor{verde}{\underline{RS2.1:}}} El Hastag debe empezar por "\#". \textbf{RF:} RF2.1. \textbf{RD(s):} RDW2.1. \\
\hspace*{1cm} \textit{Descripción:} Si el hashtag no había sido creado y no empieza por "\#", no se inserta el nuevo hashtag y se devuelve un error. \\[6pt]

\textbf{\textcolor{verde}{\underline{RS2.2:}}} El Hashtag puede ya existir o ser uno nuevo. \textbf{RF:} RF2.1. \textbf{RD(s):} RDW2.1 y RDR2.1. \\
\hspace*{1cm} \textit{Descripción:} El hashtag introducido puede que no existiera aún en la base de datos o que si existiera. En caso de ya existir, se actualiza el contador de veces que ha aparecido ese hashtag sumándole uno; si no, se crea y se pone el contador a uno.\\[6pt]
\pagebreak

\section{Listar tendencias}
\textbf{\textcolor{rojo}{\underline{RF2.2:}}} Listar tendencias. \\[6pt]

\textbf{\textcolor{rojo}{\underline{Descripción:}}} Mostrar los 10 hashtags con mayores apariciones.
 \\[6pt]

\textbf{\textcolor{rojo}{Entrada:}} Agente externo: usuario. Acción: solicitar listado de tendencias. 
Requisito de datos de entrada: ninguno. \\[6pt]

\textbf{\textcolor{rojo}{BD:}} Requisito de datos de lectura \textcolor{azul}{RDR2.2}\\[6pt]

\textbf{\textcolor{rojo}{Salida:}} Agente externo: usuario. Acción: confirmación del resultado.
Requisitos de datos de salida: \textcolor{azul}{RDS2.2}. \\[6pt] 

\textcolor{azul}{\textbf{RDR2.2:}} Datos para ver cuáles son los hashtags con más apariciones. \\
\hspace*{1cm} Los mismos que RDW2.1\\

\textcolor{azul}{\textbf{RDS2.1:}} Datos para listar tendencias. \\
\hspace*{1cm} Los mismos que RDW2.1\\

\pagebreak

\section{Asignar una categoría a un hashtag}
\textbf{\textcolor{rojo}{\underline{RF2.3:}}} Asignar una categoría a un hashtag. \\[6pt]

\textbf{\textcolor{rojo}{\underline{Descripción:}}} Asignar una categoría a un hashtag ya existente para poder clasificarlos.\\[6pt]

\textbf{\textcolor{rojo}{Entrada:}} Agente externo: administrador. Acción: solicitar actualizar la categoría de un hashtag. 
Requisito de datos de entrada \textcolor{azul}{RDE2.3} \\[6pt]

\textbf{\textcolor{rojo}{BD:}} Requisito de datos de lectura \textcolor{azul}{RDR2.3} y de escritura \textcolor{azul}{RDW2.3}\\[6pt]

\textbf{\textcolor{rojo}{Salida:}} Ninguna\\[6pt] 

\textcolor{azul}{\textbf{RDE2.3:}} El hashtag y la categoría que se le quiere asignar. \\
\hspace*{1cm} Hashtag: Cadena de caracteres (32)\\[12pt]
\hspace*{1cm} Categoría: Cadena de caracteres (32)\\[12pt]

\textcolor{azul}{\textbf{RDR2.3:}} Hashtag que hay que comprobar si ya existe en el sistema. \\
\hspace*{1cm} El mismo que \textcolor{azul}{RDR2.1}\\

\textcolor{azul}{\textbf{RDW2.3:}} Categoría a asignar. \\
\hspace*{1cm} Categoría: Cadena de caracteres (32)\\[12pt]

\textbf{\textcolor{verde}{Restricciones Semánticas:}} \\[6pt]
\textbf{\textcolor{verde}{\underline{RS2.3:}}} El hashtag debe existir previamente en el sistema. \textbf{RF:} RF2.3. \textbf{RD(s):} RDW2.3 y RDR2.3. \\
\hspace*{1cm} \textit{Descripción:} Si el hashtag no se encuentra almacenado en el sistema, no se crea ni se le asigna la categoría, sino que se devuelve un error. \\[6pt]

\pagebreak

\section{Mostrar categoría ordenada}
\textbf{\textcolor{rojo}{\underline{RF2.4:}}} Mostrar categoría ordenada. \\[6pt]

\textbf{\textcolor{rojo}{\underline{Descripción:}}} Listar las tendencias de una determinada categoría ordenadas de mayor a menos número de apariciones.\\[6pt]

\textbf{\textcolor{rojo}{Entrada:}} Agente externo: usuario. Acción: solicitar mostrar una categoría. 
Requisito de datos de entrada \textcolor{azul}{RDE2.4} \\[6pt]

\textbf{\textcolor{rojo}{BD:}} Requisito de datos de lectura \textcolor{azul}{RDR2.4}\\[6pt]

\textbf{\textcolor{rojo}{Salida:}} Agente externo: usuario. Acción: confirmación del resultado.
Requisitos de datos de salida: \textcolor{azul}{RDS2.4}. \\[6pt] 

\textcolor{azul}{\textbf{RDE2.4:}} La categoría que se quiere mostrar. \\
\hspace*{1cm} El mismo que \textcolor{azul}{RDW2.3}\\

\textcolor{azul}{\textbf{RDR2.4:}} Datos necesarios para ver si un hashtag pertenece a la categoría dada y ordenarlos por su número de apariciones. \\
\hspace*{1cm} Los mismos que \textcolor{azul}{RDE2.3} +  \textcolor{azul}{RDE2.4}\\

\textcolor{azul}{\textbf{RDS2.4:}} Datos para listar la categoría. \\
\hspace*{1cm} Los mismos que  \textcolor{azul}{RDR2.4}\\

\textbf{\textcolor{verde}{Restricciones Semánticas:}} \\[6pt]
\textbf{\textcolor{verde}{\underline{RS2.4:}}} La categoría debe existir previamente en el sistema. \textbf{RF:} RF2.4. \textbf{RD(s):} RDR2.4. \\
\hspace*{1cm} \textit{Descripción:} Si la categoría no se encuentra almacenada en el sistema se devuelve un error. \\[6pt]

\pagebreak

\section{Eliminar tendencia}
\textbf{\textcolor{rojo}{\underline{RF2.5:}}} Eliminar tendencia. \\[6pt]

\textbf{\textcolor{rojo}{\underline{Descripción:}}} Hacer que un hashtag deje de ser tendencia poniendo su contador de apariciones a cero.\\[6pt]

\textbf{\textcolor{rojo}{Entrada:}} Agente externo: administrador. Acción: solicitar la eliminación de una tendencia. 
Requisito de datos de entrada \textcolor{azul}{RDE2.5} \\[6pt]

\textbf{\textcolor{rojo}{BD:}} Requisito de datos de lectura \textcolor{azul}{RDR2.5} y de escritura \textcolor{azul}{RDW2.5}\\[6pt]

\textbf{\textcolor{rojo}{Salida:}} Ninguna

\textcolor{azul}{\textbf{RDE2.5:}} Hashtag que se quiere eliminar. \\
\hspace*{1cm} El mismo que \textcolor{azul}{RDE2.1}\\

\textcolor{azul}{\textbf{RDW2.5:}} Datos que hay que modificar para hacer que el hashtag deje de ser tendencia. \\
\hspace*{1cm} Contador: Entero\\[12pt]

\textcolor{azul}{\textbf{RDR2.5:}} Hashtag que se quiere eliminar. \\
\hspace*{1cm} El mismo que \textcolor{azul}{RDE2.1}\\

\textbf{\textcolor{verde}{Restricciones Semánticas:}} \\[6pt]
\textbf{\textcolor{verde}{\underline{RS2.5:}}} El hashtag debe existir previamente en el sistema. \textbf{RF:} RF2.5. \textbf{RD(s):} RDW2.5 y RDR2.5. \\
\hspace*{1cm} \textit{Descripción:} Si el hashtag no se encuentra almacenado en el sistema se devuelve un error. \\[6pt]

\pagebreak

